% Clase 8 --- Cuadrado de cargas: potencial en el centro (§3.2).
% Las cuatro cargas están a la MISMA distancia d = A/raiz(2) del centro, así
% que el potencial sólo pide sumar las cargas: no hay que proyectar nada.
\begin{center}
\begin{tikzpicture}[scale=1]

  % cuadrado
  \draw[figaux] (-1.5,-1.5) rectangle (1.5,1.5);

  % cargas en los vértices
  \node[figpos] at (-1.5,1.5) {};  \node[figlbl, anchor=south east] at (-1.55,1.55) {$Q_1$};
  \node[figpos] at (1.5,1.5) {};   \node[figlbl, anchor=south west] at (1.55,1.55) {$Q_2$};
  \node[figneg] at (1.5,-1.5) {};  \node[figlbl, anchor=north west] at (1.55,-1.55) {$Q_3$};
  \node[figpos] at (-1.5,-1.5) {}; \node[figlbl, anchor=north east] at (-1.55,-1.55) {$Q_4$};

  % centro
  \node[figneutra, minimum size=5pt] at (0,0) {};
  \node[figlbl, anchor=north west] at (0.1,-0.08) {$O$};

  % medias diagonales, todas iguales
  \foreach \x/\y in {-1.5/1.5, 1.5/-1.5, -1.5/-1.5} {
    \draw[figgray, line width=0.6pt] (0,0) -- (\x,\y);
  }
  \draw[figvecaux] (0,0) -- (1.32,1.32);
  \node[figlblaux, anchor=north west] at (0.72,0.66) {$d$};

  % lado
  \draw[figvecaux] (0,-2.15) -- (-1.5,-2.15);
  \draw[figvecaux] (0,-2.15) -- (1.5,-2.15);
  \node[figlblaux, anchor=south] at (0,-2.1) {$A$};

  % lectura
  \node[figlbl, figred, anchor=west, align=left] at (2.6,0.35)
    {$d = \dfrac{A}{\sqrt{2}}$ \ para las cuatro};
  \node[figlblaux, anchor=west, align=left] at (2.6,-0.75)
    {no hay que proyectar:\\ sólo sumar las cargas};

\end{tikzpicture}\\[0.5em]
{\small\itshape Como el centro equidista de los cuatro vértices, el
$1/R_{0i}$ sale de factor común y el potencial queda proporcional a
$\sum_i Q_i$ ---la carga \textbf{total}, con su signo. Nótese el contraste: el
\textbf{campo} en $O$ no es cero ni sale de una suma tan simple, porque cada
aporte apunta en otra dirección.}
\end{center}
